\documentclass[12pt]{article}
\usepackage[utf8]{inputenc}

%%% I added a few packages. This next one allows you to set the margins.
\usepackage[top=.5in,bottom=1in,right=.75in,left=.75in]{geometry}
\usepackage{multicol}
%\usepackage{newspaper}
\setlength{\columnsep}{1cm}
\usepackage{blindtext}
\usepackage{varwidth}

%%% This one does boxes in a really nice way.
\usepackage{tcolorbox}
\usepackage{mdframed}
\usepackage{minibox}
\usepackage{graphicx}
\graphicspath{ {./images/} }
\usepackage{epstopdf}
  \epstopdfDeclareGraphicsRule{.pdf}{png}{.png}{convert #1 \OutputFile}
  \DeclareGraphicsExtensions{.png,.pdf}
\usepackage{hyperref}
\hypersetup{
    colorlinks=true,
    linkcolor=blue,
    filecolor=magenta,      
    urlcolor=cyan,
    pdftitle={MSCSMess},
    pdfpagemode=FullScreen,
    }
\urlstyle{same}
 \usepackage{fix-cm}
 \usepackage{soul}
 \usepackage{graphicx}
 \usepackage{multicol}


\newcommand{\ds}{\displaystyle}
\newcommand{\ZZ}{\mathbb{Z}}
\newcommand{\QQ}{\mathbb{Q}}
\newcommand{\RR}{\mathbb{R}}
\newcommand{\CC}{\mathbb{C}}
\newcommand{\FF}{\mathbb{F}}
\newcommand{\xx}{\mathbf{x}}
\newcommand{\uu}{\mathbf{u}}
\newcommand{\vv}{\mathbf{v}}

\newcommand{\bolder}{\textbf}
 
 
\title{\fontsize{60}{70}\selectfont {\bf {MSCS} 
\includegraphics[scale=.3]{images/OleLion.png} {Mess}}}
\author{Department of Mathematics, Statistics, and Computer Science \\ St.\@ Olaf College, Northfield, MN 55057 \vspace{-2ex}} 
\date{\today\, $|$ Volume 54, No.\@ 7  \line(1,0){500}}


\begin{document}

\maketitle \vspace{-0.5cm}

Welcome back from Fall Break! We hope you had a chance to rest and enjoy the beautiful fall colors. As we dive in to the second half of the semester, we’re excited to share this week’s updates — including \textbf{course descriptions} for the upcoming J-term and Spring semester! Take some time to explore these exciting opportunities and have a great weekend!


 \begin{multicols}{2}

\begin{tcolorbox}[colback=white, colframe=black]
    \begin{center}
    \title{{\large \bolder{Women in Mathematics (AWM)}}}
    \end{center}
\end{tcolorbox}

We are excited announce the to kick off a new semester by the Association for Women in Mathematics (AWM)! Whether you are passionate about math, looking to build community, or just curious about what they do - they would love to have you involved. Please fill out this \href{https://docs.google.com/forms/d/e/1FAIpQLScYRxJUW1Kez2SKb9FY_rZgEhTlFKsD0dvfugCnlX8Vg2LVsw/viewform}{AWM Interest Form}, so they can get to know you and keep you in the loop. All are welcome - you don't have to be a math major to join. Hope to see you there!


\begin{tcolorbox}[colback=white, colframe=black]
    \begin{center}
    \title{{\large \bolder{Statistics and Data Science Club!
}}}
    \end{center}
\end{tcolorbox}

Hello everyone! The Statistics and Data Science Club is back, and they would love to have you join! Apply your data skills to real-world projects, explore mentorship opportunities, and grow your expertise beyond the classroom. Sign up for the email list \href{https://docs.google.com/forms/d/e/1FAIpQLSeV00ldRadWwtvfA8nmOIuL43iZ0MFHp6O5cmrPY5HAsxVJMQ/viewform}{here} to get updates about meetings and events. 



\begin{tcolorbox}[colback=white, colframe=black]
    \begin{center}
    \title{{\large \bolder{Problem Solving Group}}}
    \end{center}
\end{tcolorbox}
The Math Problem Solving Group meets weekly to have fun, solve math problems, eat candy, and prepare for upcoming math contests. Meetings are 11:30am-12:30pm, in RNS 207 (the Center for Interdisciplinary Research room). Everyone is welcome, regardless of prior experience in mathematics or math contests.

The first contest will be the Jerry Heuer Team Contest on the morning of Saturday, November 15. It's a fun, accessible contest for teams of up to three students. For more information, or to be added to the problem solving group email list, contact Prof. Wright at wright5@stolaf.edu.


\begin{tcolorbox}[colback=white, colframe=black]
    \begin{center}
    \title{{\large \bolder{MSCS Lightning Talks}}}
    \end{center}
\end{tcolorbox}

On Friday, October 3rd and Friday, October 24th, at 3:30pm, the MSCS department will be hosting lightning-quick research talks from our faculty. This is a great opportunity to get a glimpse at what kinds of high-level research is done by professional mathematicians, and might help inspire some questions of your own! Come to RNS 210 to check them out!



\begin{tcolorbox}[colback=white, colframe=black]
    \begin{center}
    \title{{\large \bolder{DAAD RISE Internships}}}
    \end{center}
\end{tcolorbox}

Are you interested in a summer STEM internship/fellowship program for undergraduates in Germany, requiring no German-language skills? Check out the DAAD Research Internships in Science and Engineering (RISE) website \href{https://www.daad.de/rise/en/}{at this link}. The application window is October 15th - November 30th.


 \begin{tcolorbox}[colback=white, colframe=black]
    \begin{center}
    \title{{\large \bolder{ Prospective Graduate Student Visit Day}}}
    \end{center}
\end{tcolorbox}
We’re excited to announce the "Prospective Graduate Student Visit Day" in the CS department at the University of Iowa on Friday, November 7! This is a great chance to learn more about the program and connect with faculty and current graduate students. Check out more information in the poster below!


 \begin{tcolorbox}[colback=white, colframe=black]
    \begin{center}
    \title{{\large \bolder{Math Mocktail Mixer!}}}
    \end{center}
\end{tcolorbox}
Monday, November 3, 3:00-4:30 PM, on the 6th Floor of RMS. Do any of the listed MSCS courses make you excited to learn more? Come join us for mocktails and freshly popped popcorn! Faculty and other students will be available to talk about MSCS courses as you prepare to register. This is a chance to ask any questions you may have about course topics, content, schedule, etc as you think about what you might want to take in January and Spring.

\end{multicols} \vspace{0.5cm}

\begin{center}
    {\Large \textbf{Upcoming Math Courses}}
\end{center}

With January term and Spring term registration just around the corner, the math professors have put together some advertisements for the upcoming course offerings!

\begin{multicols}{2}
\begin{tcolorbox}[colback=white, colframe=black]
    \begin{center}
    \title{{\large \bolder{SDS 282: Survey Design and Analysis}}}
    \end{center}
\end{tcolorbox}

\noindent{\textbf{Prereqs:}} SDS 172 and SDS 164. 

The Statistics and Data Science program will be offering a new elective course this J term titled Survey Design and Analysis taught by Professor Carlos Chavez. This course explores the theories and best practices for designing and conducting surveys, and statistical methods commonly used when analyzing survey data. Ask Prof. Chavez for more details!  

\begin{tcolorbox}[colback=white, colframe=black]
    \begin{center}
    \title{{\large \bolder{MATH 226: Multivariable Calculus}}}
    \end{center}
\end{tcolorbox}

\noindent{\textbf{Prereqs:}} Math 126/128 and Math 220.

This course extends important ideas of single-variable calculus (derivatives, integrals, graphs, approximation, optimization, fundamental theorems, etc.) to higher-dimensional settings. These extensions make calculus tools far more powerful in modeling the (multi-dimensional) real world. Topics include partial derivatives, multiple integrals, transformations, Jacobians, line and surface integrals, and the fundamental theorems of Green, Stokes, and Gauss.


\begin{tcolorbox}[colback=white, colframe=black]
    \begin{center}
    \title{{\large \bolder{MATH 230: Differential Equations I}}}
    \end{center}
\end{tcolorbox}

\noindent{\textbf{Prereqs:}} Math 126/128 and Math 220.

In this course, we will study Ordinary Differential Equations (ODEs), which show up in many application areas and describe the relationship between a function and its rates of change. We will study the basic theory and solution techniques for ODEs, including first-order and higher-order linear equations, systems of ODEs, and their qualitative behavior. Other methods, such as numerical methods and Laplace transforms will be introduced. We will investigate modeling real-world problems such as over-fishing in lakes, radioactive decay chains, bungee jumping, predator-prey relationships, and more! 


\begin{tcolorbox}[colback=white, colframe=black]
    \begin{center}
    \title{{\large \bolder{MATH 234: Discrete Mathematical Reasoning}}}
    \end{center}
\end{tcolorbox}

\noindent{\textbf{Prereqs:}} Math 126/128 or CSCI 221.

Want to try a different flavor of mathematics? Try Discrete Mathematical Reasoning, the study of discrete objects via mathematical reasoning! You have encountered in your mathematical education continuous functions (e.g., polynomials) that can be studied with calculus. In discrete mathematics, we will need to use new tools to understand non-continuous phenomena. The trade-off is that discrete objects can be counted and enumerating them allows us to describe real-world situations. We will also learn about the language of mathematics (logic and proof writing) and practice it across many areas, such as Set Theory, Number Theory, Combinatorics, and Graph Theory. This course is required for computer science majors, and may not be taken after MATH 244 or MATH 252.


\begin{tcolorbox}[colback=white, colframe=black]
    \begin{center}
    \title{{\large \bolder{MATH 236: Mathematics of Biology}}}
    \end{center}
\end{tcolorbox}

\noindent{\textbf{Prereqs:}} Math 126/128 and some familiarity with R.

This course is an introduction to the interdisciplinary field of mathematical biology. We will investigate a wide variety of biological systems and ask questions such as “How do tumors grow and how can we best treat them?”, “How do infectious diseases spread in a population?”, and “How can we prevent an invasive species from eradicating local plants?”.
These questions are best attacked from a variety of angles and we will employ different techniques from biological experiments (this course has an actual wet lab component!), mathematical modeling, and analysis of discrete and continuous dynamical systems to statistical analysis and computational studies using R. None of us will be experts in all of these and the intention of this course is to bring together students with many different backgrounds in Math and/or Biology to collaborate on projects - which is just what research looks like in the “real world”!  This course counts toward neuroscience, public health studies, and mathematical biology concentrations.


\begin{tcolorbox}[colback=white, colframe=black]
    \begin{center}
    \title{{\large \bolder{MATH 242: Modern Computational Mathematics}}}
    \end{center}
\end{tcolorbox}

\noindent{\textbf{Prereqs:}} Math 220.

This course is about using computational tools, namely computer algebra systems and programming languages, to explore and discover mathematical concepts. We will study topics such as algorithms for computing pi, the Fibonacci sequence and its generalizations, iterated functions and chaos, prime numbers, and probabilistic simulation. If you take this course, you will not only improve your coding skills, but also gain practice using computational experiments to provide evidence for claims, stimulate questions, and formulate precise conjectures. 


\begin{tcolorbox}[colback=white, colframe=black]
    \begin{center}
    \title{{\large \bolder{MATH 244: Real Analysis I}}}
    \end{center}
\end{tcolorbox}

\noindent{\textbf{Prereqs:}} Math 126/128 and Math 220.

The goal of Real Analysis is to master, through proof writing, the concepts in calculus (e.g., limit, continuity, derivatives, and integrals) necessary to verify such important results as the Fundamental Theorem of Calculus and the Mean Value Theorem. Students are introduced to the precision and sophistication of theoretical mathematics, and in doing so, are prepared to encounter applications involving optimization, Taylor and Fourier series, and the myriad other applications of calculus.


\begin{tcolorbox}[colback=white, colframe=black]
    \begin{center}
    \title{{\large \bolder{MATH 252: Abstract Algebra I}}}
    \end{center}
\end{tcolorbox}

\noindent{\textbf{Prereqs:}} Math 220.

Do you want more structure in your life? At least \textit{abstractly}? Come learn all about algebraic structures, their operations, their sub-structures, nice maps between them, and even quotients! For example, in linear algebra you might remember looking at vector spaces, vector addition and scalar multiplication, subspaces, linear transformations, and in particular, maps with interesting kernels. Want to know more about which rules, or \textit{axioms}, make these objects so nice? The amazing thing is that algebraic structures are ubiquitous in mathematics: once you learn about what groups and rings are, you will find them \textit{everywhere}!


\begin{tcolorbox}[colback=white, colframe=black]
    \begin{center}
    \title{{\large \bolder{MATH 262: Probability Theory}}}
    \end{center}
\end{tcolorbox}

\noindent{\textbf{Prereqs:}} Math 126/128.

``Chance governs all.’’ -  John Milton.

A great evolution of human thinking is the realization that there will always be many things unknown to us.  Things which are not certain.  Probability Theory is the mathematician’s attempt at taming uncertainty through a judicious use of measurement.  Using probability theory frees us from the worry of certainty and allows us to analyze the secret mechanisms of the world without fear.   A little bit of calculus (MATH 126/128) will suffice to take the first steps!  Some experience with mathematical software (e.g., Mathematica) is nice, but not strictly necessary.


\begin{tcolorbox}[colback=white, colframe=black]
    \begin{center}
    \title{{\large \bolder{MATH 266: Operations Research}}}
    \end{center}
\end{tcolorbox}

\noindent{\textbf{Prereqs:}} Math 126/128 and Math 220; Math 226 and/or Math 262 recommended.

Operations research, in general, seeks to answer the following question: How can we find the best solution to a problem when faced with some constraints? For example: how does an airline optimize the routes of their airplanes, given they have a limited fleet? Or: how can a company maximize their profits when faced with production constraints? And: how can a specialist determine the safest, most effective radiotherapy treatment plan given a patient’s particular cancer location and radiation beam limitations?   We’ll delve into techniques to solve problems like these (and many more).


\begin{tcolorbox}[colback=white, colframe=black]
    \begin{center}
    \title{{\large \bolder{MATH 340: Complex Analysis}}}
    \end{center}
\end{tcolorbox}

\noindent{\textbf{Prereqs:}} Math 220 and Math 226 or 244.

Complex analysis treats the calculus of complex-valued functions of a complex variable. Familiar words and ideas from ordinary calculus (limit, derivative, integral, maximum and minimum, infinite series) reappear in the complex setting. In this setting, many of these ideas feel more at home, and a deeper beauty behind them becomes apparent! Topics include complex mappings, derivatives, and integrals; applications focus especially on the physical sciences. This year, the special topic of complex dynamics will be introduced, with particular attention paid to the world-famous Mandelbrot set fractal!


\begin{tcolorbox}[colback=white, colframe=black]
    \begin{center}
    \title{{\large \bolder{MATH 364: Combinatorics}}}
    \end{center}
\end{tcolorbox}

\noindent{\textbf{Prereqs:}} Math 244 or 252.

Combinatorics is the branch of mathematics that explores clever techniques behind counting, arranging, and analyzing discrete structures. In this course, we’ll learn how to bring logic and creativity together to think combinatorially.  This means our solutions will not involve pushing algebraic symbols and variables around a page, but rather on reasoning through how a collection of objects can be counted or described in multiple different ways.  These creative explanations reveal connections between common mathematical structures and allow for deeper overall understanding.  Let's just say: you'll never see the Fibonacci numbers the same way ever again!


\begin{tcolorbox}[colback=white, colframe=black]
    \begin{center}
    \title{{\large \bolder{MATH 384: Stochastic Processes}}}
    \end{center}
\end{tcolorbox}

\noindent{\textbf{Prereqs:}} Math 262 and Math 244 or 252.

%\noindent{\textbf{Note:}} This course counts as an elective for the SDS Concentration (classes of 2026/2027) and as a sequence with MATH 262 for the Math Major.

This seminar is an introduction to stochastic processes, i.e. random objects that have entire trajectories as their outcomes. These processes are used to model a variety of real life phenomena, such as tracking the state of a randomly evolving population, a moving particle, or stock indices over time.  We will look at discrete-time Markov chains first, then study Poisson processes (briefly mentioned in Probability Theory), and finally combine the two into continuous-time Markov processes. We will answer questions such as "Which percentage of time is spent in a certain state?", "What is the probability (and expected time) to end up in certain states?", or more applied "What is the probability of a population to survive long-term?" and "How long do I expect to wait in line at my local grocery store?"  We will first build up the basic theory and tools together and then work on projects that dive deeper into applications or additional theoretic aspects that are of interest to individual students. 


\begin{tcolorbox}[colback=white, colframe=black]
    \begin{center}
    \title{{\large \bolder{MATH 396: Directed Undergraduate Research in Operator Algebras}}}
    \end{center}
\end{tcolorbox}

\noindent{\textbf{Prereqs:}} Math 244 and Math 252, admission by application.

An operator algebra is an algebraic structure that is closed with respect to a certain topology. This course will focus on operator algebras that arise from $k$-dimensional directed graphs, called $k$-graphs. In fact, we will focus more on the $k$-graphs and less on the algebras, though they will always be lurking in the background. A $k$-graph is a network of points (called vertices) that are connected by arrows (called edges) together with a set of equivalence classes on pairs of edges. We will investigate moves that can be done to a $k$-graph that leave the associated operator algebra unchanged. Since a $k$-graph can be represented by matrices, a $k$-graph move can also be realized through multiplication by specific matrices. We may also attempt to identify matrices that implement some existing $k$-graph moves.


\begin{tcolorbox}[colback=white, colframe=black]
    \begin{center}
    \title{{\large \bolder{EDUC 350: Teaching of Mathematics}}}
    \end{center}
\end{tcolorbox}

Are you interested in teaching mathematics? Whether you’re considering becoming a college math professor, a middle or high school math teacher (but aren’t pursuing math education at St. Olaf), or you’re simply curious about the student-centered approach to teaching and learning math, EDUC 350 is for you! This course aims to help you understand what it means to provide authentic mathematical experiences, why they matter, and how to develop the pedagogical skills needed to offer those experiences to your future students.

If you’re interested in taking EDUC 350 and/or would like to know more about the course, please contact Prof. Matsuura \\(matsuura@stolaf.edu).




    
\end{multicols}

\vfill

\begin{center}
\includegraphics[scale=0.385]{images/Math Mocktail Mixer.png}
\end{center}

\begin{center}
\includegraphics[scale=0.385]{images/SDS282.png}
\end{center}

\begin{center}
\includegraphics[scale=0.8]{images/VisitingStudentSchedule-2025.png}
\end{center}

\begin{center}
\includegraphics[scale=0.7]{images/Tea Time Poster.png}
\end{center}





\noindent{To submit an article, event, or anything else for publication in the Mess, email messeditor@stolaf.edu. Visit \href{https://wp.stolaf.edu/mscs/mscs-mess/}{http://wp.stolaf.edu/mscs/mscs-mess/} for a digital archive of previous MSCS Mess issues.}\\
	
\hfill Nima Dahir, Editor
	
\hfill Kyle Salois, Faculty Advisor


 

\end{document}
